\documentclass{exam}
\usepackage{graphicx} % Required for inserting images
\usepackage{fancyhdr}
\usepackage{amssymb}
\usepackage{amsmath}
\usepackage{amsthm}
\usepackage{mathtools}
\usepackage[most]{tcolorbox}
\usepackage[dvipsnames]{xcolor}
\usepackage{blindtext}
\usepackage{upgreek}
\setlength{\parindent}{0pt}
% if you ever want to enter images use the following code%
%\begin{center}%
    %\includegraphics[width=0.8\textwidth]{Screenshot (22).png}%
%\end{center}%

\newtheorem{thm}{Theorem}[section]
\newtheorem{lem}{Lemma}[section]
\newtheorem{defn}{Definition}[section]
\newtheorem{cor}{Corollary}[section]
\newtheorem{axm}{Axiom}[section]

\fancyhf{}
\fancyhead[R]{   Omkar Tasgaonkar}
\renewcommand\headrulewidth{0.5pt}
\pagestyle{fancy}
\author{Omkar Tasgaonkar}
\setlength{\headsep}{0.5in}
\fancyfoot[C]{(\thepage)}
\renewcommand\footrulewidth{0.5pt}

\fancypagestyle{firstpage}{
    \fancyhf{}
    \fancyhead[C]{\textbf{Topology: Bases and Sequential Characterizations}}
    \fancyhead[R]{Omkar Tasgaonkar}
    \renewcommand\headrulewidth{0.5pt}
    \fancyfoot[C]{(\thepage)}
    \renewcommand\footrulewidth{0.5pt}
}



\begin{document}
\thispagestyle{firstpage}

\section*{Global Bases}
Given some global set $X$ (and for further purposes, we will use a metric space $(X, \varrho)$), we define the basis as a smaller subcollection than the topology that "generates" the topology. It is similar to a framework, from which we can "derive" the elements in a topology.

\begin{defn}A global basis (or base) for a topology (notated by $\mathcal{B}$) must adhere to the following axioms:
\begin{enumerate}
  \item $\forall x \in X, \exists B \in \mathcal{B}: x \in B$
  \item $\forall B_1, B_2 \in \mathcal{B}, \exists B_3 \in \mathcal{B}: B_3 \subseteq B_1 \cap B_2$
\end{enumerate}
\end{defn}
For the remainder of this topic, we will refer to a global basis as a basis and add on the word "local" to differentiate a local basis.

\medskip
We see that part one of the definition contains part of the definition of being in the union of something. Given some arbitrary collection $A$:
\begin{equation*}
x \in \bigcup A \iff \exists B \in A: x \in B
\end{equation*}
From part one of Defn 0.1, we get the following implication
\begin{equation*}
  \forall x: x \in X \implies \exists B \in \mathcal{B}: x \in B
\end{equation*}
The reverse implication is by the construction of the basis. Since its a collection of subsets of $X$, if $\exists B \in \mathcal{B}: x \in B$, then $x \in X$. We now see that $\mathcal{B}$ is a cover for $X$:
\begin{defn} An cover for a set $X$ is a collection of subsets $U_\alpha$:
\begin{equation*}
  X \subseteq \bigcup U_\alpha
\end{equation*}
\end{defn}
Note the one way inclusion. In the case of the basis, it acts as an cover such that $X = \bigcup \mathcal{B}$. This however does not give us the notion of what sets are open vs closed (in the topological sense), so the basis generates a topology.

\begin{defn} We define a set to be "open" as the following:
  $O \in \uptau \iff \forall x \in O, \exists B \in \mathcal{B}: x \in B$
\end{defn}
Here we do not have an existing topology, rather we have a basis for the set (which follows the axioms in Defn 0.1). We now create a topology using the above definition, allowing us to distinguish which sets are open and which sets are closed. (We can verify this is a valid topology by checking the axioms of a topology and using Defn 0.1)\\

Since trivially $\forall B \in \mathcal{B}$, all elements in $B$ belong in the basis element itself, each basis element is open. Thus our cover for $X$ can be called an \textbf{open cover}. The set $X$ is trivially open then (by definition of topology; the open cover agrees with this).\\

Now, $\mathcal{B}$ is more than just an open cover for all of $X$. We see the same structure for open subsets of $X$:
\begin{lem}
\begin{equation*}
  \forall O \in \uptau, \exists B_O \subseteq \mathcal{B}: O = \bigcup B_O
\end{equation*}
where $B_O$ is a collection of basis elements.
\end{lem}
We will skip over the proof of this lemma, but in essence, the goal is to show inclusions both ways. What this tells us that the topology generated by the basis elements $\uptau$ is the same topology as the basis $\mathcal{B}$ and all its arbitrary unions and finite intersections.
\pagebreak
\section*{Local Bases}
From the concept of global bases appears the concept of a local basis (or a basis around a specific point), which gives us a look into the topology around a specific point. We will take a look at the specific definition of a local basis and then explain the logic to construct it:
\begin{defn} A local basis $\mathcal{B}_x$ (in a purely topological sense) is defined as the following:
\begin{equation*}
  \forall O \in \uptau: x \in O \implies \exists B_x \in \mathcal{B}_x: x \in B_x \subseteq O
\end{equation*}
\end{defn}
Notice how this requires a predefined topology (we will not explore the case in which we do not have a defined topology as that dips more heavily into the field of topology). We will introduce one more preliminary concept:
\begin{defn}
A filter base $\mathcal{B}$ is "the basis of a filter" with the following properties:
\begin{enumerate}
  \item $\varnothing \notin \mathcal{B}$
  \item $\forall B_1, B_2 \in \mathcal{B}, \exists B_3 \in \mathcal{B}: B_3 \subseteq B_1 \cap B_2$
\end{enumerate}
\end{defn}
We bring up this definition because our local basis $\mathcal{B}_x$ acts as a filter base for the neighborhood filter of $x$ or $\mathcal{N}_x$. Definition 0.4 is equivalent, but that requires diving into filters and directed sets, which we will not do for this set of notes.\\

Before we pivot to sequential characterizations, we discussed the notion of topological closure as the following:
\begin{equation*}
  \forall A \subseteq X: x \in \bar{A} \iff \forall B_x \in \mathcal{B}_x: B_x \cap A \neq \varnothing
\end{equation*}
and that we can write the closure as the disjoint union:
\begin{equation*}
  \bar{A} = \{\text{limit points}\} \cup \{\text{isolated points}\} = \{\text{adherent points}\}
\end{equation*}
Observe how in our definition, all points in the local basis intersect with the subset $A$. What if we are able to use "directed" or nested sets to approach a point! If we can approach through countable amount of those, then we have a sequence!
\begin{defn}
A space $X$ is first countable if we can find a countable local basis for each point $x \in X$.
\end{defn}
The benefit of being able to do this is that for any open neighborhood of $x$, there will always exist one of the countable basis elements that is a subset. We will see why this helps shortly.\\
\begin{thm}
There exists a nested local basis for all first countable spaces.
\end{thm}
\begin{proof}
Consider $B_1$ and $B_2$ (basis elements from the local basis) such that $B_2 \subseteq B_1$. Then by Defn 0.5 Property 2, $\exists B_3: B_3 \subseteq B_1 \cap B_2 = B_2$.\\

Repeating this process inductively, we can find $B_n \subseteq \bigcap_{k=1}^{n-1}B_k = B_{n-1}$, thus getting the nested property $B_{n+1} \subseteq B_n$. The fact that these are countable comes from first countability.
\end{proof}
\begin{cor}
Metric spaces are first countable.
\end{cor}
We can show this by using the local basis $\mathcal{B}_x \coloneq \left\{B\left(x, \frac{1}{n+1}\right)\right\}$ for all $n \in \mathbb{N}$. Then by definition of an open set, there exists $B(x, r) \subseteq O$. So given any such open ball $B(x, r)$ for any $r \in \mathbb{R}$ and $r > 0$, then $\exists n \in \mathbb{N}: r(n+1) > 1$ (Archimedian Property).\\

Manipulating the inequality, we get that:
\begin{equation*}
  0 < \frac{1}{n+1} < r
\end{equation*}
Thus we can find a ball $B\left(x, \frac{1}{n+1}\right) \subseteq B(x, r) \subseteq O$, satisfying Defn 0.4.\\

\begin{lem}
If $X$ is first-countable, then every convergent net has a countable convergent subnet (sequence).
\end{lem}
\begin{proof}
Nets (which utilize directed open neighborhoods to establish convergence) turn into sequences. Consider the following, let $O$ be our directed set of open neighborhoods:
\begin{equation*}
  \omega: O \xrightarrow{} A
\end{equation*}
is a net that maps the directed sets to points $x_o$ creating the net $(x_o)_{o \in O}$. For this net to converge:
\begin{equation*}
  \forall N \in \mathcal{N}_x \exists \tilde{o} \in O, \forall o \subseteq \tilde{o}: x_o \in N
\end{equation*}
We can use the local bases to generate a subnet or a subsequence we call:
\begin{equation*}
  s: \mathcal{B}_x \xrightarrow{} A
\end{equation*}
For this to converge, it must satisfy the following:
\begin{equation*}
  \forall N \in \mathcal{N}_x, \exists b_0 \in \mathcal{B}_x, \forall b \subseteq b_0: x_b \in N
\end{equation*}
By the definition of local base, we can always find some $b_0 \subseteq N$. Combining that with the existence of our nested construction, the subnet converges. Since the net $\omega$ converges to $x$, then the convergent subnet $s$ must also converge to $x$, thus we have a convergent sequence to $x$.\\

Note that we DID NOT use metrics here, yet we defined sequences, showing us that the concept of a sequence is independent of a metric space. However, certain properties like Cauchy cannot be defined without a metric.
\end{proof}

Equipped with the idea of adherent points and topological closure and how it is related to local bases, we can combine first countability to get sequential characterizations.

\section*{Sequential Characterizations (Closedness)}
\begin{thm}
A subset $A \subset X$ is sequentially closed if:
\begin{equation*}
  \forall \{x_n\}_{n \in \mathbb{N}} \in A^\mathbb{N}: x_n \xrightarrow{} x \implies x \in A
\end{equation*}
\end{thm}
Essentially all convergent $A$-valued sequences converge to a point in $A$, it is then sequentially closed. Recall topologically closed means $A$ contains all its adherent points. Now if we have a first countable space, we can see these notions are equivalent.\\

Using first countability, we can reduce the necessity of a point being adherent to nonempty intersections of a countable set of nested local basis elements and the set $A$. This can be characterized by a sequence, which as we defined was a convergent subnet.\\

\begin{thm}
In metric spaces, topologically closed is equivalent to sequentially closed:
\begin{equation*}
  A \ \text{closed} \iff \forall \{x_n\}_{n \in \mathbb{N}} \in A^\mathbb{N}: x_n \xrightarrow{} x \implies x \in A
\end{equation*}
\end{thm}
\begin{proof}
This statement is true in general for first countable spaces, but we will handle metric spaces specifically.\\

$\implies:$ If $A$ is topologically closed, then adherent points belong in the subset. We can rewrite sequence convergence as:
\begin{equation*}
  \forall \epsilon > 0, \exists n_0 \geqslant 0, \forall n \geqslant n_0: x_n \in B(x, \epsilon)
\end{equation*}
Here we see the local basis of nested balls pop up (it is however not a countable local basis since the sequence is predefined). Since $x_n \in A$, then $B(x, r) \cap A \neq \varnothing$, making $x$ an adherent point. Thus $x \in \bar{A} = A$.\\

$\impliedby:$ The converse is tricker. We proceed by proof by contrapositive:
\begin{equation*}
  A \ \text{not closed} \implies \exists \{x_n\}_{n \in \mathbb{N}} \in A^\mathbb{N}: x_n \xrightarrow{} x \land x \notin A
\end{equation*}
If $A$ is not closed, then it does not contain all of its adherent points. Therefore, we must construct a sequence whose limit is an adherent point that is not in $A$. We take the help of all of the groundwork lemmas and theorems mentioned in the previous section.\\

Consider an $x \notin A$ but in $\bar{A}$, then we can construct a convergent net. Using first countability and nested bases or in short Lemma 0.2, we can select $x_n \in B\left(x, \frac{1}{n+1}\right)$ for all $n \in \mathbb{N}$.\\

We obtain a sequence such that $x_n \xrightarrow{} x \land x \notin A \land x \in \bar{A}$. 
\end{proof}

\section*{Sequential Characterizations (Continuity)}
We shift our attention to solely metric spaces. Consider two metric spaces $(X, \varrho_X)$ and $(Y, \varrho_Y)$:
\begin{defn}
A function $f: X \xrightarrow{} Y$ is (pointwise continuous) at $x_0 \in X$ if:
\begin{equation*}
  \forall \epsilon > 0, \exists \delta > 0, \forall x \in X: \varrho_X(x, x_0) < \delta \implies \varrho_Y(f(x), f(x_0)) < \epsilon
\end{equation*}
\end{defn}
What pointwise continuity, to which going forward we will refer to solely as continuity, is the notion of "can I draw this function without lifting my pencil". In essence, it checks that there are no discontinuities or jumps at a certain point in the function's domain. Given some output range, there can always be inputs arbitrarily close to the point (we will see this is not true for discontinuities like jumps).

\begin{lem}
An alternative formulization for such continuity is:
\begin{equation*}
  \forall \epsilon > 0, \exists \delta > 0, \forall x \in X: f(B_X(x_0, \delta)) \subseteq B_Y(f(x_0), \epsilon)
\end{equation*}
\end{lem}
\begin{proof}
Deriving equivalence is a matter of applying properties of subsets and images. We can rephrase Defn 0.7 in terms of open balls:
\begin{equation*}
  \forall \epsilon > 0, \exists \delta > 0, \forall x \in X: x \in B_X(x_0, \delta) \implies f(x) \in B_Y(f(x_0), \epsilon)
\end{equation*}
By definition of image:
\begin{equation*}
  \forall \epsilon > 0, \exists \delta > 0, \forall x \in X: f(x) \in f(B_X(x_0, \delta)) \implies f(x) \in B_Y(f(x_0), \epsilon)
\end{equation*}
We then apply the definition of subset to get:
\begin{equation*}
  \forall \epsilon >0, \exists \delta > 0, \forall x \in X: f(B_X(x_0, \delta)) \subseteq B_Y(f(x_0), \epsilon)
\end{equation*}
\end{proof}

Now how do we prove something is not continuous at a point? Let's look at the negation:
\begin{equation*}
  \exists \epsilon > 0, \forall \delta > 0, \exists x \in X: \varrho_X(x, x_0) < \delta \land \varrho_Y(f(x), f(x_0)) \geqslant \epsilon
\end{equation*}
Now it can be quite difficult to find such an epsilon and such an $x \in X$ that satisfies the statement (it would be a matter of guess and check rather than construction), so let us look at another sequential characterization, this time of continuity.\\

We will first provide a proof for the following theorem, then link it to the ideas of nets and local bases:
\begin{thm}
\begin{equation*}
  f \ \text{continuous} \iff \forall \{x_n\}_{n \in \mathbb{N}} \in X^\mathbb{N}, \forall x \in X: x_n \xrightarrow{} x \implies f(x_n) \xrightarrow{} f(x)
\end{equation*}
\end{thm}
\begin{proof}
  $\implies$: We will use the fact that $f$ is continuous. Now consider an $X$-valued sequence converging to $x$:
  \begin{equation*}
    \forall \delta > 0, \exists n_0 \geqslant 0, \forall n \geqslant n_0: \varrho_X(x_n, x) < \delta
  \end{equation*}
  We have that $\forall \epsilon > 0, \exists \delta > 0, \forall x \in X: \varrho_X(x, x_0) < \delta \implies \varrho_Y(f(x), f(x_0)) < \epsilon$. So for each $\epsilon$, we are guaranteed some $\delta$ for which we are guaranteed some index $n_0$, so we get:
  \begin{equation*}
    \forall \epsilon > 0, \exists n_0 \geqslant 0, \forall n \geqslant n_0: \varrho_Y(f(x), f(x_n)) < \epsilon
  \end{equation*}
  which gives us the convergence $f(x_n) \xrightarrow{} f(x)$, proving the forward direction. For the converse we use contrapositive.\\

  $\impliedby$: If $f$ is not continuous then:
  \begin{equation*}
    \exists \epsilon > 0, \forall \delta > 0, \exists x \in X: \varrho_X(x, x_0) < \delta \land \varrho_Y(f(x), f(x_0)) \geqslant \epsilon
  \end{equation*}
  Equivalently:
  \begin{equation*}
    \exists \epsilon > 0, \forall \delta > 0, \exists x \in X: f(B_X(x, \delta)) \supset B_Y(f(x), \epsilon) \implies f(B_X(x, \delta)) \setminus B_Y(f(x), \epsilon) \neq \varnothing
  \end{equation*}
  So using this fact, we need to construct a sequence (this requires Axiom of Choice, first countability, nested local basis properties) that converges to $x$ and $f(x_n) \nrightarrow f(x)$. We use $\delta \coloneq \frac{1}{n+1}$ for all $n \in \mathbb{N}$.\\
  
  Choosing $x_n \in B_X(x, \delta)$ means $\forall n \in \mathbb{N}: \varrho_X(x_n, x) < \frac{1}{n+1}$ showing that $x_n \xrightarrow{} x$. Then by applying the function, we get that $f(x_n) \in f(B_X(x, \delta)) \setminus B_Y(f(x), \epsilon)$. Thus $f(x_n)$ can be a convergent sequence, but it will never converge to $f(x_n)$ as any term is strictly an $\epsilon$ or more away from it.
\end{proof}
Now we discuss the connections to nets and previously introduced ideas we touched upon on the proof. We require the following lemma:
\begin{lem}
If $A \subseteq B$, then $f(A) \subseteq f(B)$
\end{lem}
\begin{proof}
The image of a set $A$ is $f(A) \coloneq \{y \in Y: (\exists x: f(x) = y)\}$. Thus:
\begin{equation*}
  y \in f(A) \implies \exists x \in A: f(x) = y
\end{equation*}
Since $x \in B$, then $f(x) = y \in f(B)$, thus $f(A) \subseteq f(B)$. 
\end{proof}
Now consider $\delta' < \delta$, then $f(B_X(x, \delta')) \subseteq f(B_X(x, \delta))$. Because we have a countable local basis, we defined $\forall n \in \mathbb{N}: \delta_n = \frac{1}{n+1}$. By Lemma 0.2, we have a converging countable net $\omega: B_X(x, \delta_n) \xrightarrow{} x_n$.\\

The image of this net $f(\omega)$ could also be a converging net. We know that by Lemma 0.4, the nested structure of the directed set remains, so we need to show that if $\omega \xrightarrow{} x$, then $f(\omega) \xrightarrow{} f(x)$ is equivalent to continuity. While we can do a rigorous proof just like we did with sequences, the idea of the proof would be the same. Instead we will explore the intuitive side.\\

For the image of the net to converge, we must satisfy the definition explored in Lemma 0.2 for any neighborhood. Now we can talk about arbitrary sets, but since we are talking in the context of metric spaces, by Lemma 0.1, any open neighborhood is the union of basis elements, so we can simply work with the basis elements themselves.\\

Take any $B_Y(f(x), \epsilon)$. If the image net were to converge, then $\exists \delta > 0: f(B_X(x, \delta)) \subseteq B_Y(f(x), \epsilon)$. Applying first countability, there exists a specific such $\delta_n$. Notice how we have found the definition of pointwise continuity. Convergence of the net is guaranteed by Lemma 0.4 as the images of the nested bases are also nested. 
\end{document}